\documentclass[10pt,twocolumn,letterpaper]{article}

% -----------------------------------------------------------------------
% SPARKS Workshop @ CVPR 2026
% cvpr.sty already provides: times, xcolor, graphicx, amsmath, amssymb,
%   booktabs, natbib, enumitem, cleveref — do NOT reload them.
% -----------------------------------------------------------------------
\usepackage[pagebackref]{hyperref}  % must come BEFORE cvpr
\usepackage[review]{cvpr}           % review = line numbers + CVPR header
\usepackage{multirow}               % not in cvpr.sty
\usepackage{microtype}              % not in cvpr.sty

% Required by cvpr.sty for header/watermark (use any placeholder values)
\def\confName{SPARKS@CVPR}
\def\paperID{XXXX}

% -----------------------------------------------------------------------
\begin{document}

\title{Event-Driven Spacecraft Pose Estimation via\\
CNN Direct Regression from Voxel Grids}

\author{Santhosh S R S \quad Kashyap Hegde Kota \quad
        Pruthvi Janga \quad Bharatesh Chakravarthi \\
        \small{Arizona State University} \\
        \texttt{\{ssrs1, kkota3, pjanga, bshettah\}@asu.edu}
}

\maketitle

% -----------------------------------------------------------------------
\begin{abstract}
We present a deep learning pipeline for 6-DoF spacecraft pose estimation
from event-camera data, developed for the SPARKS Challenge at CVPR 2026.
Given a stream of asynchronous events from a 1280$\times$720 neuromorphic
camera, we convert each 100\,ms event window into a compact 3-channel
exponential-decay voxel grid and regress pose directly using a
ResNet-50 backbone with dual translation and rotation heads.
Rotation is supervised with geodesic distance on SO(3) rather
than Euclidean quaternion loss.  We compare three architectural
variants—a frame-by-frame CNN, a CNN+GRU temporal model, and a
heatmap-based keypoint regressor with PnP—on the SPADES dataset
(300 synthetic training sequences, 31 real test sequences).
Our best submission achieves a composite pose error of \textbf{2.063}
(translation 0.4105, orientation 1.6526\,rad) on the final leaderboard.
Analysis reveals that the synthetic-to-real domain gap, driven
primarily by a 100$\times$ timestamp-unit mismatch between training
and test data, is the dominant performance bottleneck, and that
single-frame prediction outperforms temporal modelling at the
10\,Hz pose sampling rate used in this dataset.
\end{abstract}

% -----------------------------------------------------------------------
\section{Introduction}
\label{sec:intro}

Accurate knowledge of a spacecraft's pose—position and orientation
relative to a camera—is essential for autonomous rendezvous, proximity
operations, and on-orbit servicing.  Traditional frame-based cameras
suffer from motion blur and limited dynamic range in the high-contrast
conditions of space, motivating the use of \emph{event cameras}
\cite{gallego2020event}.  Event cameras report per-pixel brightness
changes asynchronously with microsecond temporal resolution, offering
high dynamic range, low latency, and low power consumption.  These
properties are well suited to the rapid rotational motions and extreme
illumination conditions encountered in orbit.

The SPARKS challenge \cite{sparks2026} provides the SPADES dataset:
300 photorealistic synthetic training sequences rendered with event-camera
simulation, paired with 31 real event-camera recordings of a Proba-2
spacecraft mockup.  The task is to predict the 6-DoF pose (translation
$\mathbf{t}\in\mathbb{R}^{3}$ and unit quaternion
$\mathbf{q}\in\mathbb{H}$) at 10\,Hz from the raw event stream.

We develop an end-to-end CNN pipeline and explore three directions:
\begin{enumerate}
  \item \textbf{Frame-by-frame regression}: A ResNet-50 backbone with
        dual regression heads trained on 3-channel event voxel grids.
  \item \textbf{Temporal modelling}: A CNN+GRU network that aggregates
        context over 10 consecutive frames.
  \item \textbf{Keypoint + PnP}: A heatmap decoder regressing 2D
        projections of known 3D spacecraft keypoints, followed by
        iterative PnP.
\end{enumerate}
Our main findings are: (i) frame-by-frame CNN prediction outperforms
temporal aggregation at 10\,Hz; (ii) rotation error is substantially
harder to close than translation error across the synthetic-to-real gap;
and (iii) a 100$\times$ timestamp-unit mismatch between training and test
data is a key source of domain shift that future work should address explicitly.

% -----------------------------------------------------------------------
\section{Related Work}
\label{sec:related}

\paragraph{Event camera representations.}
Raw events cannot be directly ingested by conventional CNNs.
Zhu~\etal~\cite{zhu2019unsupervised} introduced the \emph{voxel grid},
which distributes events into temporal bins via bilinear interpolation,
preserving both temporal structure and polarity.  Alternative
representations include time surfaces \cite{lagorce2017hots},
event frames \cite{rebecq2017real}, and learned event-to-frame
conversions \cite{rebecq2019events}.  We use a signed voxel grid and
a three-channel exponential-decay variant that is directly
compatible with ImageNet-pretrained RGB backbones.

\paragraph{Spacecraft pose estimation.}
Classical model-based trackers (e.g., IPPE, EPnP) require textured
CAD models and struggle under event-camera noise.  Deep learning
approaches for frame-based cameras include direct regression
\cite{kendall2015posenet} and keypoint-based methods
\cite{park2019pix2pose,hu2019segmentation}.  The SPEED
\cite{sharma2019pose} and SPEED+ \cite{park2022speed} datasets
catalysed learning-based methods for spaceborne pose estimation,
but use standard cameras.  SPADES \cite{spades2024} is, to our
knowledge, the first large-scale dataset pairing event cameras
with 6-DoF spacecraft pose labels.

\paragraph{Loss functions for rotation.}
Naive $\ell_2$ loss on quaternion components is non-geometric and
breaks under the double-cover of SO(3).  Geodesic distance,
$2\arccos(|\hat{\mathbf{q}}\cdot\mathbf{q}|)$, is the natural metric
on the unit sphere and has been shown to improve convergence for
rotation regression \cite{mahendran2017so3,peretroukhin2020smooth}.
We adopt geodesic rotation loss throughout.

\paragraph{Domain adaptation.}
The DANN framework \cite{ganin2016domain} uses a gradient-reversal
layer to learn domain-invariant features.  We incorporate DANN into
our CNN+GRU model to encourage synthetic-to-real transfer.

% -----------------------------------------------------------------------
\section{Method}
\label{sec:method}

\begin{figure*}[t]
  \centering
  \includegraphics[width=\linewidth]{figure/architecture.png}
  \caption{Overview of our end-to-end pipeline.  Raw asynchronous events
  from the event camera are accumulated into a 3-channel exponential-decay
  voxel grid (or 5-channel signed voxel grid), which is passed through a
  ResNet-50 backbone.  The resulting feature vector feeds dual regression
  heads producing a 3-D translation vector and a unit quaternion for
  orientation.  Figure generated with AI assistance.}
  \label{fig:architecture}
\end{figure*}

\subsection{Event Representation}
\label{sec:repr}

Each pose timestamp $\tau$ is associated with a window of events
$\mathcal{E}_\tau = \{(x_i, y_i, p_i, t_i)\}$ spanning
$[\tau - 100\text{\,ms},\,\tau]$, where $x_i,y_i$ are pixel
coordinates, $p_i\in\{-1,+1\}$ is polarity, and $t_i$ is timestamp.

\paragraph{3-channel exponential decay (our primary representation).}
We divide the 100\,ms window into three equal sub-windows
$[w_0, w_1, w_2]$ of $\approx$33\,ms each.  For sub-window $k$ ending
at $t_k^{\text{end}}$, each event contributes a signed, exponentially
decaying weight:
\begin{equation}
  w_i = p_i \cdot \exp\!\left(-\frac{t_k^{\text{end}} - t_i}{\tau_{\text{decay}}}\right),
  \quad \tau_{\text{decay}} = 30\text{\,ms}.
  \label{eq:decay}
\end{equation}
Events are accumulated into a spatial grid per sub-window and
normalised channel-wise as $f(v) = \operatorname{sgn}(v)\log(1+|v|)$.
The resulting 3-channel tensor is compatible with ImageNet-pretrained
convolutional kernels.

\paragraph{5-channel signed voxel grid (baseline).}
Following Zhu~\etal~\cite{zhu2019unsupervised}, we divide the window
into $B=5$ temporal bins of 20\,ms each and bilinearly interpolate
each event across adjacent bins:
\begin{equation}
  V_k(x, y) = \sum_i p_i \max\!\bigl(0,\, 1 - |t_i^* - k|\bigr)
               \,\delta(x{-}x_i,\, y{-}y_i),
\end{equation}
where $t_i^* = (t_i - t_{\min}) / (t_{\max} - t_{\min}) \cdot (B{-}1)$
normalises timestamps to $[0,B-1]$.

\subsection{Network Architectures}

\paragraph{DirectPoseCNN (frame-by-frame).}
\label{sec:directcnn}
A ResNet-50 \cite{he2016deep} backbone pre-trained on ImageNet1K
provides the feature extractor.  The first convolutional layer is
patched to accept $C=3$ (or $C=5$) input channels by averaging the
pretrained RGB weights and replicating them across the required number
of channels, preserving the benefit of ImageNet initialisation.

Global average pooling yields a 2048-dimensional feature vector, which
feeds into two independent regression heads:
\begin{align}
  \hat{\mathbf{t}} &= \mathbf{W}_t^{(2)}\,\text{ReLU}(\mathbf{W}_t^{(1)}\mathbf{f} + \mathbf{b}_t^{(1)}) + \mathbf{b}_t^{(2)}, \\
  \hat{\mathbf{q}} &= \frac{\mathbf{g}(\mathbf{f})}{\|\mathbf{g}(\mathbf{f})\|_2},
\end{align}
where $\mathbf{f}\in\mathbb{R}^{2048}$, both intermediate layers have
128 units with ReLU activations and dropout (rate 0.2), and
$\mathbf{g}$ maps to 4 units before L2 normalisation enforces
the unit-quaternion constraint.

\paragraph{DomainAdaptivePoseNet (CNN+GRU).}
\label{sec:gru}
To exploit inter-frame context, we extend DirectPoseCNN with a
one-layer GRU \cite{cho2014gru} of hidden size 512.
A sequence of $T=10$ consecutive event frames is embedded independently
by the shared ResNet-50 backbone (the $(B,T,C,H,W)$ tensor is reshaped
to $(BT,C,H,W)$ for a single forward pass, then reshaped back to
$(B,T,2048)$).  The GRU processes the sequence and its final hidden
state feeds the dual regression heads.

Domain adaptation is implemented via a gradient-reversal
layer \cite{ganin2016domain} inserted before a binary domain
classifier $\text{FC}(512\to64\to2)$.  The reversal coefficient
$\lambda$ is annealed from 0 to 1 following a sigmoid schedule
over training epochs, encouraging the backbone to learn
domain-invariant features.

\paragraph{KeypointHeatmapNet (keypoint-based, future work).}
\label{sec:kp}
Rather than regressing pose directly, this variant regresses 2D
pixel coordinates of $K=14$ known 3D keypoints on the Proba-2 satellite,
and recovers pose via PnP.
Following the SimpleBaseline architecture \cite{xiao2018simple},
we use the ResNet-50 backbone \emph{without} global average pooling
(preserving spatial resolution) and attach a three-stage upsampling
decoder:
\begin{align*}
  &\text{ConvTranspose}(2048{\to}256,\;k{=}4,\;s{=}2)^{\times 3} \\
  &\quad\longrightarrow \text{Conv}(256{\to}K,\;k{=}1).
\end{align*}
Each decoder stage is followed by batch normalisation and ReLU.
The output $(B, K, H/4, W/4)$ heatmap is converted to sub-pixel
keypoint coordinates via differentiable soft-argmax with temperature
$\beta=100$.  Gaussian target heatmaps with $\sigma=2$\,px are
generated from projected ground-truth 3D keypoints.

\subsection{3D Keypoint Definitions}
\label{sec:keypoints}

We define $K=14$ keypoints on the Proba-2 satellite in its body frame
(metres):
\begin{itemize}
  \item \textbf{8 bounding-box corners}: $[\pm0.80,\;\pm0.30,\;\pm0.52]$
        — span the full 3D extent for maximum PnP conditioning.
  \item \textbf{6 face centres}: $[\pm0.80,0,0]$,
        $[0,\pm0.30,0]$, $[0,0,\pm0.52]$
        — ensure at least one visible keypoint under most viewpoints.
\end{itemize}
Camera intrinsics: $f_x = f_y = 1258.6$, $c_x = 640$, $c_y = 360$
(no distortion).

\subsection{PnP Pose Solver}
\label{sec:pnp}

Given predicted 2D keypoints $\{\hat{\mathbf{u}}_k\}$ and the known
3D coordinates $\{\mathbf{P}_k\}$, we solve for pose using
\texttt{cv2.solvePnPRansac} with the EPnP initialiser,
reprojection threshold $\epsilon=8$\,px, and minimum 4 inliers.
The resulting rotation vector is converted to a unit quaternion
via \texttt{scipy.spatial.transform.Rotation}.  If PnP fails
(fewer than 4 inlier correspondences), we fall back to the last
successfully predicted pose.

\subsection{Loss Functions}
\label{sec:loss}

\paragraph{Composite pose loss.}
For the direct regression models we minimise:
\begin{equation}
  \mathcal{L} = \lambda_t \cdot \|\hat{\mathbf{t}} - \mathbf{t}\|_2^2
              + \lambda_r \cdot 2\arccos\!\left(|\hat{\mathbf{q}} \cdot \mathbf{q}|\right),
  \label{eq:loss}
\end{equation}
where the rotation term is the geodesic distance on SO(3).  The
absolute value of the quaternion dot product handles the
double-cover of SO(3).  We ablate $\lambda_t$ and $\lambda_r$
in \cref{sec:ablation}.

\paragraph{Heatmap loss.}
For KeypointHeatmapNet:
\begin{equation}
  \mathcal{L}_{\text{hm}} = \frac{1}{N_{\text{vis}}}
  \sum_{k=1}^{K} \mathbf{v}_k \cdot
  \|\hat{H}_k - H_k^*\|_F^2,
\end{equation}
where $\mathbf{v}_k\in\{0,1\}$ is the per-keypoint visibility mask,
$H_k^*$ is the Gaussian target, and $N_{\text{vis}}$ is the count of
visible keypoints.

\subsection{Training Strategy}
\label{sec:training}

\paragraph{Two-phase curriculum.}
We train in two phases.
\emph{Phase 1} uses a 25\% data subset (75 sequences, 20 epochs) for
rapid architecture validation; the pass criterion is val.\ translation
$<10\%$ and rotation $<40°$.
\emph{Phase 2} fine-tunes on the full 100\% dataset (270 training,
30 validation sequences) for up to 100 additional epochs, resuming
from the Phase-1 best checkpoint.

\paragraph{Optimiser and schedule.}
We use AdamW with weight decay $10^{-4}$, initial learning rate
$5\times10^{-6}$, gradient clipping at $\|\cdot\|_{\max}=1.0$,
and \texttt{ReduceLROnPlateau} (factor 0.5, patience 5 epochs).
Training uses AMP (automatic mixed precision) with gradient scaling.
Early stopping triggers after 10 epochs without improvement in
validation loss.

\paragraph{Event-camera augmentations.}
Standard image augmentations (crop, flip) do not transfer meaningfully
to event voxel grids.  We design augmentations that simulate
event-camera failure modes:
\begin{enumerate}
  \item \textbf{RandomIntensityScale}: multiply all voxel values
        by $s \sim \mathcal{U}(0.8, 1.2)$.
  \item \textbf{RandomFrameDropout}: zero out an entire temporal
        channel with probability 0.1 (simulates event-rate drops).
  \item \textbf{SaltPepperNoise}: randomly set 1–2\% of spatial
        locations to $\pm 1$ (hot/dead pixels), applied with
        probability 0.7.
  \item \textbf{RandomErasing}: zero a random rectangle covering
        10–30\% of the frame area, with probability 0.5
        (simulates packet loss).
\end{enumerate}
No augmentations are applied at validation or inference time.

% -----------------------------------------------------------------------
\section{Experiments}
\label{sec:experiments}

\subsection{Dataset and Evaluation}
\label{sec:dataset}

The SPADES dataset \cite{spades2024} contains 300 synthetic training
sequences and 31 real test sequences captured with a 1280$\times$720
Prophesee event camera observing a Proba-2 spacecraft mockup in a
controlled environment.  Approximately 179\,000 training frames are
available after preprocessing at the 10\,Hz pose rate.
We hold out 10\% of training sequences (30 sequences) for validation.

The competition metric is the composite pose error:
\begin{equation}
  E = E_t + E_r, \quad
  E_t = \frac{\|\hat{\mathbf{t}} - \mathbf{t}\|}{\|\mathbf{t}\|}, \quad
  E_r = 2\arccos(|\hat{\mathbf{q}} \cdot \mathbf{q}|)\;\text{[rad]},
\end{equation}
evaluated over all 31 test sequences.

\paragraph{Timestamp scaling.}
Synthetic H5 files store timestamps in 100\,$\mu$s units, while real
test recordings use 1\,$\mu$s units.  This 100$\times$ mismatch
compresses the temporal distribution of events within voxel bins by a
factor of 100 at test time, constituting a severe domain shift.
All inference runs use \texttt{--test-timestamp-scale~1.0} to apply the
real-data unit convention.

\subsection{Implementation Details}

Preprocessing converts raw H5 event files into voxel grids stored as
HDF5 at 256$\times$448 (Phase 1) and 1280$\times$720 (Phase 2)
resolution, requiring approximately 1.4\,MB and 11\,MB per frame
respectively.  Training runs on a single NVIDIA A100-SXM4-80GB GPU.
At batch size 16 and full 100\% resolution, one epoch over 179\,K
frames takes approximately 2.5 hours.

\subsection{Main Results}
\label{sec:results}

\cref{tab:main} summarises all submitted results alongside the SPADES
dataset paper baselines \cite{spades2024}.

\begin{table}[h]
\centering
\small
\setlength{\tabcolsep}{4pt}
\caption{Test-set results on SPADES.  $E_t$: relative translation
error; $E_r$: geodesic rotation error (rad); $E$: composite score.
Lower is better.  SPADES baselines are the Direct and Hybrid approaches
reported in \cite{spades2024} on the same real test sequences.}
\label{tab:main}
\begin{tabular}{lccc}
\toprule
Method & $E_t$ & $E_r$ (rad) & $E$ \\
\midrule
SPADES Direct \cite{spades2024}  & 0.0513 & 1.416 & 1.47 \\
SPADES Hybrid \cite{spades2024}  & 0.0334 & 1.378 & 1.41 \\
\midrule
CNN+GRU, $T{=}10$ (ours)         & 0.6722 & 1.6113 & 2.2835 \\
DirectPoseCNN, $\lambda_t{=}2.0$, $\lambda_r{=}1.0$ (ours)
                                  & 0.4105 & 1.6526 & \textbf{2.063} \\
\bottomrule
\end{tabular}
\end{table}

\paragraph{Frame-by-frame beats temporal modelling.}
The CNN+GRU model scored 2.2835 ($E_t=0.6722$, $E_r=1.6113$\,rad)
vs.\ 2.063 ($E_t=0.4105$, $E_r=1.6526$\,rad) for DirectPoseCNN.
Translation error is 64\% higher under the temporal model, while
rotation error is marginally better (1.6113 vs.\ 1.6526\,rad),
suggesting the GRU captures orientation context but degrades
translation predictions.  At a 10\,Hz sampling rate with 100\,ms
voxel windows, consecutive frames share approximately 90\% of their
event window, providing minimal new temporal information while
increasing optimisation difficulty.

\subsection{Loss Weight Ablation}
\label{sec:ablation}

We ablate the loss weights $(\lambda_t, \lambda_r)$ while holding all
other hyperparameters fixed, resuming from the same Phase-1 checkpoint
to ensure a fair comparison.

\begin{table}[h]
\centering
\small
\setlength{\tabcolsep}{5pt}
\caption{Validation metrics for different loss weight combinations.
All runs resume from the same epoch-29 checkpoint.}
\label{tab:ablation}
\begin{tabular}{ccrrc}
\toprule
$\lambda_t$ & $\lambda_r$ & Val $E_t$ (\%) & Val $E_r$ (°) & Note \\
\midrule
2.0 & 1.0 & 14.64 & 92.46 & Submitted (best) \\
1.5 & 2.0 & 15.01 & 91.58 & Best val rotation \\
0.0 & 1.0 & $\approx$15.2 & $\approx$92.5 & Frozen trans.\ head \\
\bottomrule
\end{tabular}
\end{table}

Increasing $\lambda_r$ relative to $\lambda_t$ yields modest
improvements in validation rotation error (−0.88° at best) but at the
cost of elevated translation error (+0.37 pp).
Freezing the translation head while training the rotation head and
backbone causes the backbone features to shift, degrading translation
despite the frozen FC head.  This confirms that the shared backbone
tightly couples both tasks.

The best rotation result ($91.58°$) was achieved at epoch 37 with
$\lambda_t=1.5$, $\lambda_r=2.0$; however, the improvement was not
submitted due to time constraints.

\subsection{Domain Gap Analysis}
\label{sec:domaingap}

\cref{tab:gap} and \cref{fig:domaingap} quantify the
validation-to-test domain gap.

\begin{table}[h]
\centering
\small
\setlength{\tabcolsep}{5pt}
\caption{Synthetic-to-real domain gap for our best DirectPoseCNN.}
\label{tab:gap}
\begin{tabular}{lrr}
\toprule
Metric & Val (synthetic) & Test (real) \\
\midrule
Translation error & 14.64\% & 41.05\% \\
Rotation error    & 92.46°  & 94.69° \\
\midrule
Amplification     & $\times$2.77 & $+$2.23° \\
\bottomrule
\end{tabular}
\end{table}

\begin{figure}[h]
  \centering
  \includegraphics[width=\linewidth]{figure/domain_gap.png}
  \caption{Synthetic-to-real domain gap for both submitted models.
  Translation error amplifies by $2.77\times$ (DirectPoseCNN) and
  $3.24\times$ (CNN+GRU) from validation to test, while rotation error
  changes by less than $0.5°$, indicating the domain shift primarily
  affects translation.}
  \label{fig:domaingap}
\end{figure}

Translation error is amplified by 2.77$\times$ from validation to
test, while rotation error increases by only 2.23°.  The large
translation amplification likely stems from the timestamp-unit mismatch
(100$\times$), which alters the density of events per voxel bin.
Because rotation depends primarily on the \emph{spatial pattern} of
events and translation depends on absolute scale, the former is more
robust to temporal rescaling.

\subsection{Leaderboard Context}
\label{sec:leaderboard}

\cref{tab:leaderboard} shows the final public leaderboard.

\begin{table}[h]
\centering
\small
\setlength{\tabcolsep}{4pt}
\caption{Final SPARKS 2026 leaderboard (top 5 and our entry).}
\label{tab:leaderboard}
\begin{tabular}{clccc}
\toprule
Rank & Team & $E_t$ & $E_r$ (rad) & $E$ \\
\midrule
1 & National Key Lab & 0.0146 & 0.0355 & \textbf{0.0501} \\
2 & NSSC\_UCAS        & 0.0230 & 0.0331 & 0.0561 \\
3 & haoranhuang       & 0.1464 & 0.5562 & 0.7026 \\
4 & pruthj            & 0.1335 & 0.9101 & 1.0436 \\
5 & csu\_nuaa\_pang   & 0.3340 & 0.7724 & 1.1064 \\
\midrule
12 & santhoshsrs (ours) & 0.4105 & 1.6526 & 2.063 \\
\bottomrule
\end{tabular}
\end{table}

The top two teams achieve near-perfect translation
($E_t \approx 0.02$), suggesting they employ model-based methods
or exploit explicit 3D spacecraft geometry rather than relying on
learned translation from data alone.

% -----------------------------------------------------------------------
\section{Discussion}
\label{sec:discussion}

\paragraph{Geodesic loss is necessary.}
Preliminary experiments with $\ell_2$ quaternion loss failed to
converge reliably (rotation errors $> 150°$).  Geodesic distance in
\cref{eq:loss} provides well-defined gradients across the entire
SO(3) manifold and handles the sign ambiguity of quaternion
representation.

\paragraph{Timestamp normalisation is critical for domain transfer.}
The 100$\times$ timestamp-unit mismatch between synthetic and real data
is the dominant source of domain shift.  Normalising timestamps within
each event window to a fixed range $[0, T_{\text{window}}]$ before
constructing voxel grids would make the representation invariant to
this shift and is a straightforward improvement for future work.

\paragraph{Temporal modelling does not help at 10\,Hz.}
With 100\,ms voxel windows and 10\,Hz pose sampling, adjacent frames
are highly correlated.  Temporal modelling with a GRU introduces
additional complexity without providing new information.  A higher
pose sampling rate or shorter event windows might benefit from
temporal aggregation.

\paragraph{GlobalAvgPool destroys spatial information for keypoints.}
Our initial keypoint regressor, which applied GlobalAvgPool before the
regression head, was fundamentally limited to $\approx$160\,px
keypoint error because all spatial information was discarded.
Switching to a heatmap decoder (removing the pooling operation)
is the correct architectural choice; preliminary experiments showed
expected error of 20–40\,px.  Future work should complete the
heatmap training and evaluate the full keypoint+PnP pipeline.

\paragraph{Loss weight coupling.}
Because translation and rotation heads share the ResNet-50 backbone,
tuning $\lambda_t$ and $\lambda_r$ creates implicit trade-offs:
improving rotation can hurt translation and vice versa.  Disentangling
these tasks, for instance with separate backbone branches or
task-specific normalisation layers, could provide more independent control.

% -----------------------------------------------------------------------
\section{Conclusion}
\label{sec:conclusion}

We presented an event-camera spacecraft pose estimation pipeline
evaluated on the SPADES dataset as part of the SPARKS Challenge at
CVPR 2026.  Our best result, a composite pose error of \textbf{2.063},
is achieved by a ResNet-50 direct regression model trained with
geodesic rotation loss on 3-channel exponential-decay event voxel grids.
We found that: (i) single-frame prediction outperforms temporal
modelling at 10\,Hz; (ii) synthetic-to-real translation error is
amplified $\approx$2.8$\times$ at test time due to a 100$\times$
timestamp-unit mismatch; and (iii) loss weight tuning for the
rotation head is constrained by the shared backbone.
Future work will focus on timestamp-invariant voxel construction,
a completed heatmap keypoint regressor with PnP, and explicit domain
normalisation to reduce the synthetic-to-real gap.

% -----------------------------------------------------------------------
{\small
\bibliographystyle{ieeenat_fullname}
\bibliography{references}
}

\end{document}
